\documentclass[../../../uem_mei_q.tex]{subfiles}

\begin{document}
    次の空欄\:\myboxb{エ}\,，\,\myboxb{オ}\:に当てはまるものを指定された解答群の中から選び，%
    解答用紙の所定の欄の番号をマークせよ．なお，解答群から同じものを2回以上選んでもよい．%
    それ以外の空欄には，当てはまる0から9までの数字を解答用紙の所定の欄にマークせよ．%
    ただし，$\log$は自然対数で，$i$は虚数単位である．

    \begin{enumerate}
        \item 関数 $\displaystyle%
            f(t)=t\sqrt{t^2+1}+\log\left(t+\sqrt{t^2+1}\right)$ を微分すると \\%
            $f'(t)=\:\myboxb{ア}\:%
            \sqrt{\rule{0pt}{13pt}t\:^{\myboxbs{イ}}+\:\myboxb{ウ}\,}$ となる．
        \item $t$を実数とする．複素数 $z=(1+ti)^2$ を考える．$x$，$y$を実数として%
            $z=x+yi$ とおくと，\\
            \hspace{3\zw}%
            $\displaystyle%
                x=\:\myboxb{エ}\,,\quad y=\:\myboxb{オ}\qquad \cdots\cdots\:(*)%
            $\\
            と表すことができる．

            $t$が実数全体を動くとき，座標平面上で$(*)$によって媒介変数表示される曲線を$C$とおく．%
            $C$の $0\leqq t\leqq1$ の部分の長さは，\\
            \hspace{3\zw}%
            $\displaystyle%
                \sqrt{\,\myboxb{カ}\,}+%
                \log\left(\:\myboxb{キ}\:+\sqrt{\,\myboxb{ク}\,}\,\right)%
            $\\
            となる．

            $(*)$の2式から媒介変数$t$を消去すると，$C$は\\
            \hspace{3\zw}%
            $\displaystyle%
                x=\:\myboxb{ケ}\:-\myfraca{1}{\myboxb{コ}\,}y^2%
            $\\
            によって定まる放物線であることがわかる．$C$の焦点の座標は%
            $\left(\:\myboxb{サ}\,,\;\myboxb{シ}\:\right)$で，準線は $x=\:\myboxb{ス}$ である．%
            $C$と3つの直線 $y=0$，$y=2$，$x=\myboxa{ス}$ で囲まれた図形を$x$軸のまわりに1回転してできる%
            立体の体積は $\myboxb{セ}\:\pi$ となる．

            \vspace{.5\baselineskip}
            \textsf{エ，オの解答群}
            
            {\renewcommand{\arraystretch}{1.3}%
                \vspace{-.1\baselineskip}%
                \hspace{\zw}%
                \begin{tabularx}{\dimexpr\linewidth-2\zw}{@{}LLLL@{}}
                    ⓪\; $1+t$ & ①\; $1-t$ & ②\; $t-1$ & ③\; $t$ \\
                    ④\; $-t$ & ⑤\; $2t$ & ⑥\; $-2t$ & ⑦\; $1+t^2$ \\
                    ⑧\; $1-t^2$ & ⑨\; $t^2-1$
                \end{tabularx}%
            }

    \end{enumerate}
\end{document}
